% problem-000096 2 分光光度分析
\section*{2. 分光光度分析}
\textit{下列为分问。}

\subsection*{2-1}
等物质的量的⼆氧化碳可和钛配合物Cp TiMe （下图）在戊烷溶液中光照(350 nm)条件下发⽣“插⼊”反应。画出该反应产物的结构式。

（图：）

\subsection*{2-2}
⼆氧化碳“插⼊”到过渡⾦属配合物中的M—N键中⽣成相应的配合物。⼀个例⼦是配合物 $\boldsymbol \eta ^ { 5 } \cdot$ CpRu(NH_{2})dcpe $\mathrm { ( d c p e = C y _ { 2 } P C H _ { 2 } C H _ { 2 } P C y _ { 2 } , C y }$ : 环已基，⻅下图，配体dcpe省略）和⼆氧化碳在THF（四氢呋喃）中的反应。画出该反应的关键中间体和产物的结构式。

（图：）

\subsection*{2-3}
⼆氧化碳可“插⼊”到双⾦属配合物内的M—M′键得以固定。有研究表明，Fe-Zr双⾦属配合物$\mathrm { C p ( C O ) _ { 2 } F e - Z r ( C l ) C p _ { 2 } }$ （下图）中插⼊等物质的量的⼆氧化碳的反应过程经过⼀四元环过渡态进⽽形成产物。画出该过渡态和产物的结构式。

（图：）

\subsection*{2-4}
⽤于⼆氧化碳催化转化的钳型铱配合物(PNPyP)IrH X（X =Me, OH; Me为甲基）的结构如下图所⽰。研究表明，该结构中 $\mathrm { I r { - } H _ { a } }$ 键的键⻓较 $\mathrm { I r { - } H _ { b } }$ 键为短。

（图：）

\subsection*{2-4}
-1 该配合物可以和等物质的量的⼆氧化碳发⽣“插⼊”反应。画出相应产物的结构式。

\subsection*{2-4}
-2 配合物(PNPyP)IrH_{2}X中配体X为Me时记为A，X为OH时记为B。A和B哪个更易发⽣上述⼆氧化碳“插⼊”反应？

\subsection*{2-5}
⼆氧化碳和如下图所⽰的化合物 $\left( \mathrm { A r { \mathrm { - } } 2 } , 6 { \mathrm { - } } ^ { i } \mathrm { P r } _ { 2 } \mathrm { C } _ { 6 } \mathrm { H } _ { 3 } \right)$ 在甲苯中发⽣1:1加成反应；较之反应物，产物内的磷-磷键⻓变短⽽磷-镓键⻓变⻓。画出该产物的结构式。

（图：）
